% Synthesis and Further Directions  (source §22 + new: a synthesis "three views"
% figure, and pointers onward -- j-invariant/moduli, isogenies/torsion, curves
% over Q and F_p, modular forms and modularity)
\chapter{Synthesis and Further Directions}
\label{ch:synthesis}

We have reached the destination announced in \cref{ch:intro}. This closing
chapter recollects the argument as a single chain, then sketches the directions
the subject opens onto --- the moduli of elliptic curves, isogenies and torsion,
the arithmetic over \(\Q\) and finite fields, and the link to modular forms ---
each of which begins exactly where this book ends.

\section{The Complete Chain}
\label{sec:syn-chain}

The proof of \cref{thm:main} is the following chain, each link a chapter.

\begin{enumerate}
\item A lattice \(\Lambda \subset \C\) gives a compact genus-one Riemann surface
\(T = \C/\Lambda\) (\cref{ch:lattices,ch:torus_group_surface}).
\item \(T\) carries the intrinsic abelian group law \([u]+[v]=[u+v]\)
(\cref{prop:tg-welldef}).
\item A genus-one dimension count forces the first nonconstant function to have a
double pole --- the Weierstrass \(\wp\)
(\cref{prop:rr-counts,ch:weierstrass_wp}).
\item The functions \(1, \wp, \wp'\) span \(L(3O)\) and define a map
\(\Phi : T \to \CP^2\) (\cref{sec:emb-l3o}).
\item The differential equation \(\wp'^2 = 4\wp^3 - g_2\wp - g_3\) places the image
on the cubic \(E\) (\cref{thm:wp-ode,prop:emb-lands}).
\item The origin \(0 \in T\) maps to \(O = [0:1:0]\), and \(\Phi\) is a
biholomorphism (\cref{prop:emb-holo,thm:emb-iso}).
\item The inverse \(z \mapsto -z\) becomes \((x,y) \mapsto (x,-y)\)
(\cref{sec:emb-reflection}).
\item A line cuts \(E\) in \(P+Q+R\); since \(Z = 0\) cuts \(3O\), the divisor
\(P+Q+R-3O\) is principal (\cref{thm:lin-principal}).
\item Principal divisors on \(T\) have Abel sum zero, so \(P+Q+R = O\)
(\cref{thm:abel-law,thm:lin-collinear}).
\item Hence \(P \oplus Q = -R\): the chord--tangent law (\cref{thm:tr-main}).
\end{enumerate}

Three descriptions of one group have emerged, all canonically isomorphic
(\cref{fig:syn-threeviews}): the \emph{analytic} group \((\C/\Lambda, +)\), the
\emph{geometric} group \((E, \boxplus)\) of chord--tangent, and the
\emph{algebraic} group \(\Pic^0(E)\) of degree-zero divisor classes.

\begin{figure}[ht]
\centering
\begin{tikzpicture}[every node/.style={font=\small}]
  \node[visualnode, text width=3.0cm] (an) at (0,0)
    {\textbf{analytic}\\$(\C/\Lambda,\,+)$\\addition mod $\Lambda$};
  \node[visualnode, text width=3.0cm] (geo) at (7,0)
    {\textbf{geometric}\\$(E,\,\boxplus)$\\chord--tangent};
  \node[visualnode, text width=3.0cm] (alg) at (3.5,-2.6)
    {\textbf{algebraic}\\$\Pic^0(E)$\\divisor classes};
  \draw[arrow] (an) -- node[above, labelnode] {$\Phi$ (\cref{thm:emb-iso})} (geo);
  \draw[arrow] (geo) -- node[right, labelnode, xshift=2pt] {$P \mapsto [P]-[O]$} (alg);
  \draw[arrow] (alg) -- node[left, labelnode, xshift=-2pt] {$\bar\sigma$ (\cref{thm:div-pic-iso})} (an);
\end{tikzpicture}
\caption{Three views of one group. The Weierstrass embedding \(\Phi\) identifies
the analytic and geometric groups; the Abel--Jacobi map \(\bar\sigma\) identifies
the algebraic group \(\Pic^0(E)\) with the torus; and \(P \mapsto [P]-[O]\) closes
the triangle. All three are the same abelian group.}
\label{fig:syn-threeviews}
\end{figure}

\begin{keyidea}{Final takeaway}
The chord--tangent law is the visible projective shadow of the simple addition
law \([u]+[v]=[u+v]\) on the complex torus. The Weierstrass embedding turns
addition into line intersection because line intersections give principal
divisors, and principal divisors have zero Abel sum. The torus explains the
group; the cubic displays it; the divisor classes name it.
\end{keyidea}

\section{Further Directions}
\label{sec:syn-further}

Everything below takes \((E, O) \cong (\C/\Lambda, 0)\) as its starting point.

\subsection{The \texorpdfstring{$j$}{j}-Invariant and Moduli}
\label{sec:syn-j}

\Cref{prop:lat-normal-form} reduced the classification of complex tori to
\(\HH / \mathrm{SL}_2(\Z)\). A single number resolves it: the \emph{\(j\)-invariant}
\[
  j = 1728\,\frac{g_2^3}{g_2^3 - 27 g_3^2} = 1728\,\frac{g_2^3}{\Delta}.
\]
Two elliptic curves over \(\C\) are isomorphic iff they have the same \(j\), and
\(j : \HH / \mathrm{SL}_2(\Z) \to \C\) is a bijection: the moduli space of elliptic
curves is the \(j\)-line \(\C\). As a function of \(\tau\), \(j(\tau)\) is the
basic modular function, invariant under \(\mathrm{SL}_2(\Z)\). The lemniscatic
curve \(y^2 = 4x^3 - 4x\) of our examples has \(j = 1728\)
(\cref{ex:syn-j1728}), a curve with extra symmetry (complex multiplication by
\(\Z[i]\)).

\subsection{Isogenies and Torsion}
\label{sec:syn-isogeny}

An \emph{isogeny} is a nonconstant holomorphic group homomorphism
\(E \to E'\); analytically it is induced by \(z \mapsto \alpha z\) for some
\(\alpha \in \C\) with \(\alpha\Lambda \subseteq \Lambda'\). The most important
example is multiplication by \(n\),
\[
  [n] : E \to E, \qquad P \mapsto \underbrace{P \oplus \cdots \oplus P}_{n},
\]
a rational map of degree \(n^2\) whose kernel is the \(n\)-torsion. Over \(\C\),
\[
  E[n] = \{ P : nP = O \} = \tfrac1n\Lambda / \Lambda \cong (\Z/n\Z)^2 ,
\]
generalizing the \(E[2] \cong (\Z/2\Z)^2\) of \cref{sec:cases-torsion}. The
torsion carries the Weil pairing \(E[n] \times E[n] \to \mu_n\), a perfect
alternating pairing that pervades the arithmetic theory.

\subsection{Elliptic Curves over \texorpdfstring{$\Q$}{Q} and Finite Fields}
\label{sec:syn-arithmetic}

Because the group law is given by rational formulas (\cref{prop:coord-field}), it
runs over any field. Two cases drive modern number theory.

Over \(\Q\), the \emph{Mordell--Weil theorem} states that \(E(\Q)\) is a finitely
generated abelian group,
\[
  E(\Q) \cong \Z^r \oplus E(\Q)_{\mathrm{tors}},
\]
with \(r\) the \emph{rank}. The torsion is completely classified (Mazur); the rank
is deep, conjecturally governed by the \(L\)-function through the
Birch--Swinnerton-Dyer conjecture.

Over a finite field \(\F_p\), the group \(E(\F_p)\) is finite with
\[
  \#E(\F_p) = p + 1 - a_p, \qquad |a_p| \le 2\sqrt p
\]
(Hasse's bound). The difficulty of inverting the map \(n \mapsto nP\) in
\(E(\F_p)\) --- the elliptic curve discrete logarithm problem --- is the basis of
elliptic-curve cryptography.

\subsection{Modular Forms and Modularity}
\label{sec:syn-modular}

The Eisenstein series \(G_4, G_6\) that produced \(g_2, g_3\)
(\cref{def:wp-eisenstein}) are the first \emph{modular forms}: under
\(\Lambda \mapsto \alpha\Lambda\) they scale by \(\alpha^{-4}, \alpha^{-6}\), so as
functions of \(\tau\) they satisfy \(G_k(\tfrac{a\tau+b}{c\tau+d}) =
(c\tau+d)^k G_k(\tau)\). The discriminant \(\Delta = g_2^3 - 27g_3^2\) is a cusp
form of weight \(12\), and \(j = 1728 g_2^3/\Delta\) has weight \(0\)
(\cref{ex:syn-weight}). The arithmetic of elliptic curves over \(\Q\) is bound to
modular forms by the \emph{modularity theorem} (Wiles, Taylor--Wiles,
Breuil--Conrad--Diamond--Taylor): every elliptic curve over \(\Q\) arises from a
weight-two modular form. Its most famous consequence is Fermat's Last Theorem.

\subsection{Where to Read Next}
\label{sec:syn-reading}

Silverman's \emph{The Arithmetic of Elliptic Curves} \cite{silverman} is the
standard entry to the arithmetic theory; Silverman--Tate, \emph{Rational Points on
Elliptic Curves} \cite{silverman-tate}, is a gentler start. For the
complex-analytic and modular side, see Diamond--Shurman, \emph{A First Course in
Modular Forms} \cite{diamond-shurman}, and Serre, \emph{A Course in Arithmetic}
\cite{serre}. Each of these opens at the point reached here: the elliptic curve is
a complex torus, and its group law is addition.

\section{Summary}
\label{sec:syn-summary}

The chord--tangent law is addition on \(\C/\Lambda\) transported by
\((\wp, \wp')\), and the same group appears in three guises --- analytic,
geometric, algebraic (\cref{fig:syn-threeviews}). From this single fact the
subject branches: the \(j\)-invariant classifies the curves
(\cref{sec:syn-j}); isogenies and torsion organize the maps between them
(\cref{sec:syn-isogeny}); the rational group law carries the theory over \(\Q\)
and \(\F_p\) into number theory and cryptography (\cref{sec:syn-arithmetic}); and
the Eisenstein series tie it to modular forms and, through modularity, to Fermat
(\cref{sec:syn-modular}). The torus was the key throughout: it makes a mysterious
geometric construction into the most natural thing in the world, the addition of
points.

\section*{Exercises}
\addcontentsline{toc}{section}{Exercises}

\begin{exercise}
\label{ex:syn-j1728}
Compute the \(j\)-invariant of \(E : y^2 = 4x^3 - 4x\) (so \(g_2 = 4\),
\(g_3 = 0\)). Confirm \(\Delta = 64\) and \(j = 1728\). What is \(j\) for a curve
with \(g_2 = 0\) (so \(y^2 = 4x^3 - g_3\))?
\end{exercise}

\begin{exercise}
\label{ex:syn-torsion}
Show that over \(\C\), \(E[n] = \tfrac1n\Lambda/\Lambda \cong (\Z/n\Z)^2\), and
that multiplication by \(n\), realized as \(z \mapsto nz\) on \(\C/\Lambda\), is an
\(n^2\)-to-one map, so \([n]\) has degree \(n^2\).
\end{exercise}

\begin{exercise}
\label{ex:syn-weight}
Using \(G_k(\alpha\Lambda) = \alpha^{-k} G_k(\Lambda)\), show that under
\(\Lambda \mapsto \alpha\Lambda\) one has \(g_2 \mapsto \alpha^{-4} g_2\),
\(g_3 \mapsto \alpha^{-6} g_3\), \(\Delta \mapsto \alpha^{-12}\Delta\), and hence
\(j\) is invariant. Interpret these as the weights \(4, 6, 12, 0\).
\end{exercise}

\begin{exercise}
\label{ex:syn-hasse}
Take \(E : y^2 = x^3 + 1\) over \(\F_5\). List its affine points, add the point at
infinity, and find \(\#E(\F_5)\). Check that it satisfies the Hasse bound
\(|{\#E(\F_5)} - 6| \le 2\sqrt5\), and identify the group structure.
\end{exercise}
